\documentclass[12pt]{article}

\usepackage[margin=1in]{geometry}
\usepackage{amsmath, amssymb}
\usepackage{graphicx}
\usepackage{booktabs}
\usepackage{array}
\usepackage{caption}
\usepackage{subcaption}
\usepackage{setspace}
\usepackage{parskip}
\usepackage{hyperref}
\usepackage{natbib}
\usepackage{float}
\usepackage{multirow}
\usepackage[T1]{fontenc}
\usepackage{microtype}
\usepackage{tocloft}

\hypersetup{
    colorlinks=true,
    linkcolor=black,
    citecolor=black,
    urlcolor=black
}

% Keep the contents page compact enough to fit on one page.
\setcounter{tocdepth}{2}
\setlength{\cftbeforesecskip}{3pt}
\setlength{\cftbeforesubsecskip}{1pt}
\setlength{\cftsecnumwidth}{2.2em}
\setlength{\cftsubsecnumwidth}{3.0em}
\renewcommand{\cftsecfont}{\normalsize}
\renewcommand{\cftsubsecfont}{\small}
\renewcommand{\cftsecpagefont}{\normalsize}
\renewcommand{\cftsubsecpagefont}{\small}

\onehalfspacing

\title{\textbf{Exploring Adaptive Complexity in Momentum and Pairs Trading}\\[0.4em]
\large Empirical Finance Capstone}
\author{Akshat Koirala\\Macalester College}
\date{May 2026}

\begin{document}

\maketitle

\begin{abstract}
We compare two model-based trading strategies: one using a Bayesian Information Criterion (BIC) framework to choose the optimal time horizon each month; and the other employing a Kalman Filter to estimate the unobserved mean of the log-return spreads of foreign exchange (FX) currency pairs. Both dynamic models are compared to their static versions. Our tests run from January 2016 until December 2024. The BIC sector strategy outperforms the fixed 12-month baseline in terms of average returns ($14.12\%$ per year vs.\ $13.30\%$, respectively), as well as risk-adjusted performance (Sharpe Ratio of $0.69$ compared to $0.62$). However, the Kalman FX strategy performs slightly worse than its OLS counterpart ($1.37\%$ per year vs.\ $1.53\%$, and Sharpe Ratio of $-0.19$ vs.\ $-0.14$). These results imply that the added complexity of the two model-based strategies is worthwhile if the markets have enough regimes such that it allows for the use of multiple horizons, however, if the markets are too quickly mean reverting then this added complexity will add additional delay to the trade signals.
\end{abstract}

\newpage

\begingroup
\singlespacing
\small
\setlength{\parskip}{0pt}
\tableofcontents
\endgroup

\newpage
\onehalfspacing

%--------------------------------------------------------------
\section{Introduction}
%--------------------------------------------------------------

A continuing argument within the area of Quantitative Finance is: Does increasing complexity in the model enhance trading results, or does it simply increase the estimation errors and the costs associated with implementing them? Models with increased levels of complexity (e.g., additional parameters) place greater demands upon the cost of implementing these models due to the need to provide highly sensitive calibration. Additionally, many complex models produce delayed trade signals compared to simpler models. Therefore, rather than asking if a more complicated model can be fit to past data (i.e., in sample), we want to know if this produces better trade signals once the costs of producing these signals have been subtracted.

To address this question, I test two separate trading methodologies independently. One was a Sector Rotation Strategy. For this strategy, an investment decision is made monthly regarding which of eight U.S. Equities sectors' ETFs should be held. In addition, this decision is made by looking at each sector's recent momentum. There is a much simpler implementation (the base case) and one where the model included Bayesian Information Criterion (BIC). The BIC model selected how far back to use in determining momentum for each sector but the base case looks back twelve months for all sectors.

The second trading methodology tested in this study was Statistical Arbitrage Trading Methodology. With this methodology, an investment decision is made to purchase and/or sell currency pair FX spreads that fall outside of their estimated mean. Again, there were two versions of this strategy: one with a simple model (base case) and one that utilized a Kalman Filter (complexity). The base case used a static OLS z-score to estimate the current spread state (i.e., the "true" spread state) while the complex model utilized the Kalman Filter to estimate the current spread state.

I paired these two methodologies together intentionally. Each comparison maintained everything else equal: same universe, same entry/exit criteria, same transaction costs, and changed only the way in which the signal was generated. This allowed me to isolate the incremental effect of the complexity in the modeling methodology.

In total, the back test ran from January 2016 through December 2024. All of the testing was done over this time frame. The years 2015--16 were used as a warm up to determine optimal parameters for each of the models. The initial capital balance was \$100,000. Equity trades incurred a transaction fee of 5 bps (\$0.05) per dollar traded. FX trades had a round-trip commission fee of 3.00 pips (\$30.00 USD per pip x 2 sides = \$60.00 USD).

The remainder of this paper will follow the following structure:

Section 2 provides an overview of my Data. Section 3 outlines the Methodology for both of the trading methods examined. Section 4 outlines and displays results for both methodologies for the entire testing period. Section 5 examines results for both methodologies across different subperiods and also across different combinations of parameter values. Section 6 summarizes results, identifies potential limitations, and suggests future research opportunities.

%--------------------------------------------------------------
\section{Data}
%--------------------------------------------------------------

\subsection{Equity Data}

The sector momentum strategy uses daily adjusted close prices for eight SPDR Select Sector ETFs: XLK (Technology), XLV (Health Care), XLE (Energy), XLF (Financials), XLY (Consumer Discretionary), XLP (Consumer Staples), XLI (Industrials), and XLU (Utilities). Prices are downloaded from Yahoo Finance and SPY is included as the market benchmark. Data covers the past decade form January 2015 through December 2024; the first year is used as warmup for our parameters.

\subsection{FX Data}

The pairs trading strategy uses daily closing prices for 18 major and cross-currency pairs: EURUSD, GBPUSD, AUDUSD, NZDUSD, USDCAD, USDCHF, USDJPY, EURGBP, EURJPY, GBPJPY, AUDJPY, NZDJPY, EURCHF, GBPCHF, EURCAD, GBPCAD, AUDCAD, and AUDNZD. All prices are sourced from Yahoo Finance. Cross rates not directly quoted are computed from USD legs. These 18 instruments are arranged into 24 two-legged baskets; the basket pairing logic is described in Section 3.2.

The primary limitation of the FX dataset is that Yahoo Finance reports mid prices at the end of each day, which are not necessarily based on executable bid/offer quotes. Cross-rates can sometimes be derived from USD-based exchange rate quotations instead of being directly quoted in the market. I have added a 1.5 pips transaction cost to help address these issues; however, the true trading costs will likely be much greater due to various market volatility.

%-------------------------------------------------
\section{Methodology}
%-------------------------------------------------

\subsection{BIC-Optimized Sector Momentum}

\subsubsection{Strategy Logic}

Sector momentum strategies seek to take advantage of the observed phenomenon that industries which have performed well historically will likely do so again in the future during some horizon period. Historically, the most common method of implementing this has been to rank the sectors by their last 12-month performance and then rotate your money into the highest performing sectors on a monthly basis.

However, there exists the question if the use of a 12-month lookback is optimal across all market conditions. In rapidly rotating environments, like the 2022 rate hike environment where energy and finance were quickly separated from tech, the ability of a 12-month lookback to identify meaningful trends in the space may be inadequate. Conversely, in slower trending environments, the use of shorter time frames may introduce unnecessary volatility. BIC allows for a systematic and data driven approach to determining what the optimal lookback should be at each rebalance date without having to specify it manually.

\subsubsection{BIC Lookback Selection}

To determine the optimal lookback for each sector at the beginning of each month, an AR(1) model is estimated using the log of daily returns for each of the eight sector EFT's using one of three possible lookback windows: 63 trading days (approximately 3 months), 126 trading days (approximately 6 months), or 252 trading days (approximately 12 months). For each window, BIC is calculated based upon the formula below:

\begin{equation}
\text{BIC} = k \ln(n) - 2\ln(\hat{L})
\end{equation}

Where $k$ represents the number of estimated parameters in the AR(1) model, which is 2; the constant term and the lag coefficient, $n$ is the number of observations within the specific window being evaluated, and $\hat{L}$ represents the maximized value of the log-likelihood function. To allow for equitable evaluation of competing models regardless of the length of their respective windows, per observation BIC was employed.

For each sector, the window with the lowest BIC was designated as its corresponding lookback. The choice of lookback window was independent for each of the eight sectors; therefore, while certain sectors could utilize similar lookbacks, they did not need to.

\subsubsection{Portfolio Construction}

Following determination of the BIC-based lookback window $(L_i^{*})$ for a given sector $(i)$, the sector's momentum return $(m_{i,t})$ was determined as follows:

\begin{equation}
m_{i,t} = \frac{P_{i,t}}{P_{i, t-L^*_i}} - 1
\end{equation}

In order to construct the portfolio, each of the eight sectors had its associated momentum return measured at each monthly rebalance date. Then, using these values, each sector was ranked by its momentum return. Each month, the portfolio consisted of the top three sectors in equally weighted allocations, each allocated 33.3\%. The portfolio was rebalanced once per month and remained unchanged until the subsequent rebalance date. Additionally, a long short version of this strategy was tested wherein bottom-ranked sectors with negative momentum were also shorted in addition to the top-ranked sectors with positive momentum.

\subsubsection{Benchmark Strategy}

The benchmark strategy utilized in this research repeated exactly the same process: exact universe, identical equal-weighting methodology for selecting top-3 constituents, identical monthly rebalance schedule, identical transaction cost structure, but applied a 12-month lookback for each sector every month. Therefore, the only distinction between the BIC strategy and its benchmark counterpart was how the lookback window was determined.

%-------------------------------------------------

\subsection{Kalman-Filtered FX Statistical Arbitrage}

\subsubsection{Formation of Pair Universe and Basket Composition}

The statistical arbitrage strategy for foreign exchange utilizes the principle that economically-related currency pairs exhibit co-movement over time. As such, significant deviations from a pair's historical co-movement pattern will likely correct themselves. Twenty-four two-legged baskets are constructed from eighteen currency-pairs as follows:

\begin{enumerate}
    \item Dollar Bloc (six baskets): EUR-GBP-AUD-NZD dollar-pairs paired against each other.
    \item USD Crosses (three baskets): USDCAD, USDCHF, USDJPY paired against each other.
    \item EUR Crosses (five baskets): EURGBP-EURCHF-EURCAD-EURJPY-GBPJPY combinations.
    \item JPY Crosses (six baskets): EUR-GBP-AUD-NZD-JPY paired against each other.
    \item CAD and CHF (four baskets): Remaining AUDCAD, GBPCAD-EURCAD, and CHF cross-combinations.
\end{enumerate}

\subsubsection{Spread Calibration Through OLS}

To measure how much each basket $(A,B)$ spreads, we use ordinary least squares on the last 252 trading days' log price data. We run the OLS model to obtain the coefficients of the relationship. We calculate the hedge ratio for each basket. The hedge ratio is the number of units of Pair B needed to offset one unit of Pair A's risk.

We also calculate the spread. The spread is the difference between the actual returns and what would be expected if the relationships were perfectly linear. Since we want to compare the spread over time, we need to center it around 0. Therefore we subtract the mean of the residuals.

Finally, we perform the Augmented Dickey-Fuller test to verify that the spread is stationary before including it in our model. For example, our tests show that our AUD/USD vs. NZD/USD basket has a p-value less than $.01$ with a hedge ratio of about $.94$.

\subsubsection{Mean Reversion of Spreads (Half Life)}

Our centered spread can be thought of as a first order autoregression.

\begin{equation}
\tilde{\eta}_t = \phi \tilde{\eta}_{t-1} + \epsilon_t,
\qquad \epsilon_t \sim \text{NID}(0, \sigma^2_\epsilon).
\end{equation}

Here we have defined $\phi$ as the coefficient for mean reversion. It will be estimated using OLS regression $y = \phi x$, where $x = \tilde{\eta}_{t-1}$, and $y = \tilde{\eta}_t$. We bound $\phi$ to $[0.05, 0.99]$ so that the process is both stationary and reasonably useful. On average, we find that for our 24 baskets in the sample, our estimated values of $\phi$ are around $0.96$, i.e., they correspond to an approximate theoretical half life of just under 16 trading days, based on $t_{1/2} = \ln(2) / -\ln(\phi)$.

\subsubsection{Estimating Parameters Using Kalman Filters}

Using Kalman filters, we treat our spread as a latent variable evolving in accordance with the AR(1) transition equation and observable with some degree of noise. Our state space representation is given below.

\begin{align}
s_t &= \phi s_{t-1} + w_t, 
& w_t &\sim \text{NID}(0,Q), \\
z_t &= s_t + v_t, 
& v_t &\sim \text{NID}(0,R).
\end{align}

We choose $Q = \max(\hat{\sigma}^2_\varepsilon, 10^{-7})$, and $R = 0.5Q$. Our filter begins at the final known centered spread along with covariance $\hat{\sigma}^2_{\tilde{\eta}}$. Each day, we pass our daily observed centered spread into the filter's update method which provides us with an estimate of the hidden state $s_t$ along with an associated variance.

Thus, we define the Kalman z score as follows.

\begin{equation}
z_t^{\text{kalman}} = \frac{\hat{s}_t}{\hat{\sigma}_{\tilde{\eta}}}.
\end{equation}

%--------------------------------------------------
\section{Results}
%--------------------------------------------------

\subsection{Sector Momentum}

\begin{table}[h]
\centering
\caption{Full-Period Performance: Sector Strategies (January 2016--December 2024)}
\label{tab:sector_perf}
\begin{tabular}{lrrrrr}
\hline
\textbf{Strategy} &
\textbf{Annual Return} &
\textbf{Volatility} &
\textbf{Sharpe Ratio} &
\textbf{Maximum Drawdown} &
\textbf{Beta} \\
& (\%) & (\%) & & (\%) & \\
\hline
SPY (Benchmark) & 14.65 & 17.84 & 0.709 & -33.72 & 1.00 \\
Sector Momentum (BIC) & 14.12 & 17.58 & 0.690 & -30.64 & 0.88 \\
Sector Baseline (12 Months) & 13.30 & 18.18 & 0.622 & -30.23 & 0.90 \\
Sector Long / Short Momentum & 2.09 & 15.76 & 0.006 & -32.67 & 0.29 \\
\hline
\end{tabular}
\end{table}

Compared to the 12-month baseline, the BIC sector strategy offered an additional 82 basis points of return per annum during this time frame, a total of 14.12\% versus 13.3\%. Its sharpness also increased to 0.69 compared to 0.62 for the base line, a more significant increase considering the two methods have similar transaction costs. While maximum draw-down was marginally greater for BIC (-30.64\%) compared to the base line (-30.23\%), the differences are negligible. BIC's beta relative to SPY (.883) was only slightly below that of the base-line (.90), as expected since BIC sometimes reverses out of positions that do not correlate as well with the broader market.

Over the course of nearly nine years, both methodologies were just about fifty basis points shy of the return realized by SPY over the same time frame. BIC has not been a better performer than SPY; it simply provides an incremental improvement over a simple twelve-month base line.

\begin{figure}[h]
\centering
\includegraphics[width=0.95\textwidth]{NAV_sector_strategies.png}
\caption{\emph{Panel A}: Cumulative net asset values (NAV) for sector strategies and SPY, January 2016--December 2024; starting capital \$100,000.}
\label{fig:sector_nav}
\end{figure}

During the post-covid recovery (2020--2021) the BIC heatmap (fig.~\ref{fig:BIC_heatmap}) showed that most often the shorter three-month windows were selected especially for the energy and financial sectors. These sectors experienced large and rapid reversals. Conversely, during the 2017--2019 bull run, the twelve-month windows were almost always chosen; as would be expected: in slowly trending environments longer memory can capture more enduring momentum. The apparent variety among the colors in the heat map is, in fact, itself evidence that BIC is doing something more than superficially adapting.

\begin{figure}[h]
\centering
\includegraphics[width=0.95\textwidth]{BIC_lookback_heatmap.png}
\caption{\emph{Panel B}: BIC-selected window length (months) by sector and calendar month, 2016--2024; red = 12m, yellow = 6m, blue = 3m.}
\label{fig:BIC_heatmap}
\end{figure}

The long/short version of this methodology fared very badly at 2.09\% annualized (sharpness = 0.006). As one would expect during this time Period, there was primarily a bull-run and even when many sectors finished last they were still producing some level of positive return resulting in a consistent drag on Performance from the short side of the trade. The low beta (.285) for this methodology is consistent with much of the net-long position being hedged-away; which would explain why this methodology did little-to-nothing in terms of providing diversification benefits from the short side of trades.

\subsection{FX Statistical Arbitrage}

\begin{table}[h]
\centering
\caption{Full-Period Performance: FX Methodologies (January 2016--December 2024)}
\label{tab:fx_perf}
\begin{tabular}{lrrrrrr}
\hline
\textbf{Methodology} &
\textbf{Ann. Ret.} &
\textbf{Vol.} &
\textbf{Sharpe Ratio} &
\textbf{Max DD} &
\textbf{Beta} &
\textbf{Corr.} \\
& (\%) & (\%) & & (\%) & vs.\ SPY & vs.\ SPY \\
\hline
Kalman FX Baskets & 1.37 & 3.39 & -0.186 & -5.31 & 0.007 & 0.037 \\
FX Fixed-Beta OLS & 1.53 & 3.33 & -0.141 & -4.91 & 0.005 & 0.029 \\
SPY (Benchmark) & 14.65 & 17.84 & 0.709 & -33.72 & 1.000 & 1.000 \\
\hline
\end{tabular}
\end{table}

It may be helpful to view these FX results in conjunction with those for equities; rather than separately.

Neither methodology produced a positive Sharpe Ratio over the entire time-frame. Returns of around $1.4$--$1.5\%$ annualized for both methodologies and versus a risk-free return of $2\%$, indicate negative excess returns for both methodologies by definition.

However, what stands-out clearly in comparing these two methodologies is their risk-profiles: the annualized volatilities for both methodologies were $3.3$--$3.4\%$, roughly $1/5^{\text{th}}$ that of SPY's $17.84\%$. Similarly, while SPY suffered a maximum-draw down of $-33.72\%$; the respective methodologies suffered max draw-downs of $-5.31\%$ and $-4.91\%$. Finally, while correlation to SPY was effectively zero for both methodologies at $.037$ and $.029$ respectively, both methodologies have a higher absolute-beta value then the equivalent equity methodologies; however their beta values are still relatively small ($<.10$).

A cumulative performance chart of the FX portfolios compared to the SPY index fund from January 2016 through December 2024 has been included for reference purposes. To make it easier to see the performance of the FX portfolio, its values have been plotted along the right side of the chart rather than along the left side.

There is evidence of a difference in how these two methods perform. A comparison of the metrics produced using the OLS model and those produced by the Kalman filter method show that OLS performed better than Kalman on each of the metrics evaluated. The differences in behavior can be seen in Figure~\ref{fig:lag}. At the start of January 2016, the OLS z-score for the AUD/NZD cross currency pair passed the $+2.0$ entry threshold and triggered a buy signal immediately. The z-score of the Kalman filter, smoothed by the hidden state estimator, remained below $2.0$ at this time and therefore failed to trigger a buy signal for several days. As a result, the Kalman strategy began buying shares after their profitable window was nearly closed due to the rapid reversal of the cross currency spread.

Figure~\ref{fig:lag} does not represent a unique or isolated event. This delay is inherent to the design of the Kalman filter, which includes some degree of discounting for extreme readings as being indicative of random fluctuations rather than signals. While this may be beneficial in applications where there is a long latency associated with changes in the underlying conditions, such as inflation rates, it is detrimental in the context of foreign exchange markets where significant reversals occur within short periods of time, typically less than one week. The median value of $\phi$, calculated across all baskets, was $.956$, representing an approximate 15 day half life. In reality, however, numerous spread deviations observed in the dataset reverted in significantly shorter intervals than this, far too rapidly for the smoothing costs incurred by the Kalman filter to justify any benefits derived from reduced variance.

%--------------------------------------------------------------
\section{Robustness}
%--------------------------------------------------------------

\subsection{Performance During Sub-Periods}

In order to determine if the findings are based on a singularly favorable period during which BIC demonstrated superior returns, we divide the historical backtesting data into three separate and economically meaningful sub-periods:

\begin{table}[htbp]
\centering
\caption{Performance Results for Each Sub-Period Using Sector Strategy Data}
\label{tab:subperiod_sector}
\begin{tabular}{llrrr}
\toprule
\textbf{Sub-Period} &
\textbf{Portfolio} &
\textbf{\% Ann. Return} &
\textbf{Sharpe Ratio} &
\textbf{Max Downside Deviation (\%)} \\
\midrule
\multirow{3}{*}{Pre-COVID Period: 2016--18}
& SPY & 9.72\% & 0.595 & $-19.35\%$ \\
& Sector Momentum & 7.94\% & 0.446 & $-20.53\%$ \\
& Sector Baseline & 3.50\% & 0.104 & $-21.13\%$ \\
\midrule
\multirow{3}{*}{COVID Era: 2019--21}
& SPY & 25.93\% & 1.092 & $-33.72\%$ \\
& Sector Momentum & 23.25\% & 0.954 & $-30.64\%$ \\
& Sector Baseline & 22.74\% & 0.914 & $-30.23\%$ \\
\midrule
\multirow{3}{*}{Rate Hike Period: 2022--24}
& SPY & 8.82\% & 0.390 & $-24.50\%$ \\
& Sector Momentum & 11.80\% & 0.616 & $-14.40\%$ \\
& Sector Baseline & 13.93\% & 0.728 & $-14.40\%$ \\
\bottomrule
\end{tabular}
\end{table}

Figures showing annualized return, Sharpe ratio, and maximum drawdown for sector strategies by sub-period have also been provided.

The pre-COVID period (2016--2018) represented a general uptrend in equity prices and little dispersion among sectors. Consequently, both momentum-based sector strategies under-performed relative to SPY. BIC delivered returns of 7.94 percent compared to SPY's 9.72 percent returns. The fixed baseline strategy fared worst of all at just 3.5 percent, so BIC's ability to dynamically select among competing sector strategies was arguably more important than what appear to be headlining returns.

When considering the COVID era (2019--2021) both strategies behaved similarly to SPY. Specifically, BIC delivered returns of 23.25 percent versus SPY's 25.93 percent returns. Although BIC experienced slightly smaller max draw-downs ($-30.64$ percent for BIC vs. $-33.72$ percent for SPY) neither strategy was able to compete with a passive investment in SPY.

Finally, during the post COVID / Rate Hike period (2022--2024), results showed that BIC clearly outperformed SPY on a return basis (11.8 percent vs. 8.82 percent) and Sharpe ratio (.616 vs. .390), consistent with our hypothesis that lookback adaptation provides added value when sector leadership is changing rapidly. However, contrary to expectations, BIC was outperformed by a twelve month fixed baseline strategy in this same period (.728 Sharpe ratio vs. .616). This inconsistency represents a challenge to our simplistic interpretation that dynamic look-back adaptation produces superior returns than static look-back baselines. A possible explanation for why the dynamic strategy was outperformed by a static strategy in this particular period relates to switching between winning and losing positions prior to the twelve month signal becoming fully realized. Overall, results indicate that although BIC outperforms over the course of nine years, results do not consistently favor BIC across all sub-periods examined.

\begin{table}[htbp]
\centering
\caption{Sub-Period Performance: FX Strategies}
\label{tab:subperiod_fx}
\begin{tabular}{llrrr}
\toprule
\textbf{Period} & \textbf{Strategy} & \textbf{Ann. Ret. (\%)} & \textbf{Sharpe} & \textbf{Max DD (\%)} \\
\midrule
\multirow{2}{*}{Pre-COVID 2016--18}
 & Kalman FX        &  0.95 & $-$0.274 & $-$5.31 \\
 & FX Fixed-Beta    &  1.29 & $-$0.193 & $-$4.91 \\
\midrule
\multirow{2}{*}{COVID Era 2019--21}
 & Kalman FX        &  1.44 & $-$0.187 & $-$3.45 \\
 & FX Fixed-Beta    &  1.35 & $-$0.216 & $-$3.41 \\
\midrule
\multirow{2}{*}{Rate-Hike 2022--24}
 & Kalman FX        &  1.63 & $-$0.113 & $-$4.36 \\
 & FX Fixed-Beta    &  1.86 & $-$0.042 & $-$4.26 \\
\bottomrule
\end{tabular}
\end{table}

\begin{figure}[htbp]
    \centering
    \includegraphics[width=0.95\textwidth]{subperiod_fx.png}
    \caption{Sub-period performance for FX strategies.}
    \label{fig:subperiod_fx}
\end{figure}

The FX returns for each of the two methodologies were very comparable; each methodology generated returns within the 0.9\% to 1.9\% range for each sub-period and all methodologies had less than 5.3\% peak-to-trough drawdown. A great deal can be said about the consistency of these methodologies. These methodologies do not need a macroeconomic climate to generate fairly level real excess returns. Therefore they make good candidates for an uncorrelated overlay methodology.

\subsection{Parameter Sensitivity}

\begin{figure}[htbp]
    \centering
    \begin{subfigure}[b]{0.48\textwidth}
        \includegraphics[width=\textwidth]{robustness_heatmap.png}
        \caption{Sector: Sharpe by lookback $\times$ top-N grid}
    \end{subfigure}
    \hfill
    \begin{subfigure}[b]{0.48\textwidth}
        \includegraphics[width=\textwidth]{fx_param_sweep.png}
        \caption{FX Kalman: Sharpe by entry/exit z-score grid}
    \end{subfigure}
    \caption{Parameter sensitivity analysis for both strategies.}
    \label{fig:sensitivity}
\end{figure}

If Kalman was going to outperform OLS at any time during any sub-period it would have been during the COVID period. During that time the spread noise was high and the filtering advantage should have been greatest. It did not. The disadvantage of using lags in the Kalman methodology seems to be a characteristic of the strategy rather than something unique to the COVID era.

Sharpe Ratios for the sector methodology ranged from 0.52 to 0.73 over all combinations of fixed-lookback and top-N parameters. There were no combinations that resulted in a negative Sharpe Ratio; therefore there is some evidence that the results were not due to a single lucky choice of parameter. The highest Sharpe ratio (0.73) was achieved with a 6 month fixed-look back and choosing the three or four best sectors. However, this could only be known after the fact. The BIC method of selecting parameters adaptively produced a Sharpe ratio of 0.69 without knowing how long to wait before making the decision.

When considering the FX Kalman methodology, there was a uniform set of negative Sharpe Ratios based on varying entries into positions ($z_{entry} = 1.5$, 2.0, 2.5) and exits from positions ($z_{exit} = 0.3$, 0.5, 0.75). The worst of these configurations (entry=2.5, exit=0.5: Sharpe=$-0.102$) performed better than the OLS model. The methodology has a fundamental disadvantage because of its lag in generating returns in this market.

%--------------------------------------------------------------
\section{Discussion}
%--------------------------------------------------------------

\subsection{When Does Complexity Help?}

The sector momentum results suggest that adaptive lookback selection adds value when the market is in a regime that punishes fixed-memory models. The rate-hike period of 2022--2024 is the clearest example: rising interest rates caused sharp divergence among sectors that had previously moved together, and a strategy with a fixed 12-month backward look was slow to respond. Shorter lookbacks captured these rotations earlier, and BIC provided a mechanism to shift toward shorter windows without requiring any manual intervention.

The more general point is that BIC's value is not in being smarter than a fixed lookback on average; it is in avoiding the worst mistakes of a fixed lookback when market conditions change. This is consistent with the parameter sensitivity analysis, which shows that no single fixed lookback dominated BIC across all sub-periods.

\subsection{When Does Complexity Hurt?}

The FX results illustrate the opposite case. The Kalman filter's smoothing introduces a delay that is particularly damaging in fast mean-reverting markets. FX spread deviations, at least for the pairs and baskets examined here, tend to resolve over a few trading days. The filter's posterior estimate of the hidden state moves slowly by design; it partially attributes extreme observations to noise rather than true spread movement. In this setting, the noise being filtered out is actually part of the signal.

A useful way to think about the tradeoff is: Kalman filtering is valuable when the observation noise $R$ is large relative to the process noise $Q$, meaning the measurement is unreliable. For these FX baskets, the opposite seems to be true; the raw spread is a reasonably faithful indicator of the true mean-reversion signal, and smoothing it loses more than it gains.

\subsection{Portfolio-Level Implications}

One observation that is easy to overlook is that the FX strategies carry very low SPY correlation (0.037 and 0.029). Even though their standalone Sharpe ratios are negative, they could improve the risk-adjusted performance of a combined portfolio. The monthly returns heatmap in Figure \ref{fig:monthly_heatmap} makes this visible: the March 2020 column is dark red for SPY and the sector strategies (large losses), but the FX rows are barely off-white (near-zero losses). From a portfolio construction standpoint, a low-return, low-volatility, near-zero-correlation strategy has genuine diversification value even if its standalone Sharpe ratio is negative.

\begin{figure}[htbp]
    \centering
    \includegraphics[width=0.95\textwidth]{monthly_returns_all_heatmap.png}
    \caption{Monthly returns heatmap for all five strategies and SPY. Green = positive returns, red = losses. March 2020 column illustrates the FX strategies' diversification property.}
    \label{fig:monthly_heatmap}
\end{figure}

%--------------------------------------------------------------
\section{Limitations}
%--------------------------------------------------------------

Several limitations of this analysis deserve acknowledgment.

\textbf{In-sample BIC selection.} The BIC lookback is selected each month using all data available up to that date, which is proper walk-forward conditioning. However, the choice to test specifically the \{3, 6, 12\} month candidate set was made with knowledge of the literature, and there is no guarantee that this set is optimal out-of-sample. A true walk-forward test would require splitting the data into a design period (used to choose the candidate set) and an evaluation period.

\textbf{Data quality.} Yahoo Finance adjusted closes are a reasonable starting point, but they are not institutional-grade data. FX cross rates are synthesized from USD legs and may not reflect true executable quotes. The 5bps and 1.5-pip cost models are approximations; actual bid-ask spreads and market impact would vary with trade size and market conditions.

\textbf{No leverage or margin modeling.} The FX strategy implicitly uses leverage (holding 33\% gross exposure per open basket), but margin requirements, overnight financing costs, and forced liquidation scenarios are not modeled. In a live account, these costs would reduce performance.

\textbf{Regime sensitivity.} Both strategies performed differently across sub-periods in ways that are difficult to predict in advance. The BIC advantage materialized primarily in the 2022--2024 rate-hike period; if that period had not occurred, the full-period conclusion might be different.

%--------------------------------------------------------------
\section{Possible Extensions}
%--------------------------------------------------------------

Several natural extensions follow from this work.

\textbf{Walk-forward optimization.} A proper out-of-sample test would re-estimate the candidate lookback set (or at minimum validate the \{3, 6, 12\} choice) on a held-out period. This would provide a cleaner test of whether BIC's adaptive selection generalizes beyond the 2016--2024 sample.

\textbf{Regime filter.} Adding a macro regime overlay, for example, pausing momentum positions when the VIX exceeds a threshold, or conditioning on the slope of the yield curve, could improve robustness in crash periods. The strategy's main vulnerability is rapid market-wide drawdowns, not slow grinding bear markets.

\textbf{Intraday FX data.} The Kalman lag problem is fundamentally a frequency mismatch: the filter is designed for slower-moving signals, but the FX baskets revert quickly. Moving to hourly data would reduce the proportional lag from 2--3 days to 2--3 hours, likely eliminating the OLS advantage. This would require a proper data vendor rather than Yahoo Finance.

\textbf{Risk-parity portfolio combination.} The near-zero correlation between the FX strategies and the sector strategy is largely wasted in this analysis because they are evaluated separately. A risk-parity allocation, weighting each strategy inversely by its historical volatility, would explicitly exploit the diversification benefit and likely improve the combined portfolio Sharpe ratio meaningfully.

\textbf{Live forward test.} The only test that ultimately matters for any trading strategy is live performance. The sector and FX strategies developed here could be implemented via Interactive Brokers' API for a forward paper-trading test, which would provide genuine out-of-sample evidence.

%--------------------------------------------------------------
\section{Conclusion}
%--------------------------------------------------------------

This paper set out to test whether added model complexity improves trading outcomes, and the answer turns out to be: it depends on what kind of complexity, applied where.

For sector momentum, BIC-based adaptive lookback selection produced a modest but consistent improvement over the fixed 12-month baseline. Annualized return improved by 82 basis points and the Sharpe ratio improved from 0.622 to 0.690 over the nine-year sample. The improvement was most pronounced in the 2022--2024 rate-hike environment, where rapid sector rotation rewarded the shorter-memory selections that BIC favored. This result supports the hypothesis that complexity adds value when it tracks genuine regime changes in the underlying market, rather than fitting noise.

For FX statistical arbitrage, the Kalman filter did not improve over the fixed-beta OLS baseline on any full-period metric. The mechanism is straightforward: the filter's smoothing introduces a 2--3 day lag on entry signals, and the profit windows in fast-reverting FX spreads are short enough that this lag eliminates most of the available return. This is a useful negative result; it shows that smoothing is not universally beneficial, and that the value of the Kalman filter depends critically on the time scale of the underlying mean reversion.

The broader takeaway is that complexity should be evaluated against the market it faces, not against a simpler model in the abstract. A model that correctly identifies regime structure will outperform a fixed-parameter baseline. A model that filters signal as if it were noise will not.

%--------------------------------------------------------------
\bibliographystyle{plainnat}
\begin{thebibliography}{9}

\bibitem{quantpedia}
Quantpedia. (n.d.).
\newblock Sector momentum -- rotational system.
\newblock \textit{Quantpedia: The Encyclopedia of Quantitative Trading Strategies}.
\newblock Accessed May 7, 2026.

\bibitem{quantconnect}
QuantConnect. (n.d.).
\newblock Kalman filters and statistical arbitrage.
\newblock \textit{QuantConnect Documentation, Research Environment}.
\newblock Accessed May 7, 2026.
\newblock \url{https://www.quantconnect.com/docs/v2/research-environment/applying-research/kalman-filters-and-stat-arb}

\bibitem{jegadeesh1993}
Jegadeesh, N., \& Titman, S. (1993).
\newblock Returns to buying winners and selling losers: Implications for stock market efficiency.
\newblock \textit{The Journal of Finance}, 48(1), 65--91.

\bibitem{moskowitz1999}
Moskowitz, T.~J., \& Grinblatt, M. (1999).
\newblock Do industries explain momentum?
\newblock \textit{The Journal of Finance}, 54(4), 1249--1290.

\bibitem{kalman1960}
Kalman, R.~E. (1960).
\newblock A new approach to linear filtering and prediction problems.
\newblock \textit{Journal of Basic Engineering}, 82(1), 35--45.

\bibitem{schwarz1978}
Schwarz, G. (1978).
\newblock Estimating the dimension of a model.
\newblock \textit{The Annals of Statistics}, 6(2), 461--464.

\end{thebibliography}

\end{document}
